\begin{textsample}{POS Dim 4 – human – Score 6.00 – mg/mg\_ef\_linguagens\_ingles\_6\_p2.txt}  \label{ex:f4_pos_002}
O conhecimento linguístico ( vocabulário e gramática ), a competência \textbf{textual} ( conhecimento sobre textualidade, continuidade temática, gêneros \textbf{textuais}, tipos de \textbf{texto}, etc. ), a competência sociolinguística ( adequação da linguagem às situações de interação ) e competência estratégica ( uso consciente de estratégias para lidar com situações e contextos pouco conhecidos nas várias interações do dia-a-dia por meio da língua estrangeira, tanto na modalidade \textbf{oral} quanto na \textbf{escrita} se articulam com a finalidade de gerar sentido.
As formas gramaticais deixam de ser enfatizadas como um fim em si mesmas para serem entendidas e internalizadas como meios pelos quais é possível expressar propósitos comunicativos de acordo com o contexto das interações sociais.

% matched lemmas: escrita, oral, texto, textual
\end{textsample}
